\section{Accuracy-independent grids under quadratic growth}\label{sec:growth}

Theorem~\ref{thm:certificate} holds for any choice of grids. This section
chooses them so that, under weighted growth, every coordinate grid has
$O(\sqrt{\bar\kappa}\log n)$ nodes at every refinement stage, independently
of the accuracy reached. The growth constant is used only in the analysis.

\subsection{Graded grids and one trial}

\begin{definition}[Graded grid]\label{def:graded}
Let $B_i=[\ell_i',u_i']$ (or its integer points), a center $c_i\in B_i$ (an
integer if $i\in I_Z$), a mesh $h>0$ and a grading $0<\theta\le1/4$ be given.
Put
\begin{equation}\label{eq:step}
\sigma(t)=h+\theta t\quad(i\in I_C),\qquad
\sigma(t)=\max\{1,\lfloor h+\theta t\rfloor\}\quad(i\in I_Z).
\end{equation}
Starting from $t_0=0$, set $t_{k+1}=\min\{t_k+\sigma(t_k),u_i'-c_i\}$ until
$t_{k+1}=u_i'-c_i$, and take the nodes $c_i+t_k$; proceed in the same way
towards $\ell_i'$. These nodes form the \emph{graded grid} of $B_i$.
\end{definition}

Spacing grows linearly with the distance from the center, so a grid of radius
$R$ needs only $O(\theta^{-1}\log(1+\theta R/h))$ nodes.

\begin{lemma}[Graded grids]\label{lem:graded}
(a) $w_i(v)\le h+\theta\abs{v-c_i}$ for every node $v$.
(b) If an endpoint is at distance $R>0$ from $c_i$, the number of nodes
strictly on that side is at most $\lceil(4/\theta)\ln(1+\theta R/H)\rceil$,
where $H=h$ for $i\in I_C$ and $H=\max\{h,1\}$ for $i\in I_Z$.
(c) If $h\ge u_i'-\ell_i'$, the grid has at most three nodes.
\end{lemma}

The proof is in Appendix~\ref{app:growth}. The meshes are chosen per
coordinate so that every coordinate sees the same correction scale. For
$i\in P$ let $e_i\in\Z$ satisfy $4^{e_i-1}<L_i\le4^{e_i}$ and put
$r_i=2^{-e_i}$, so that $\frac14<L_ir_i^2\le1$. Let $E$ be the least integer
with $2^Er_i\ge s_i$ for all $i\in P$, and put $\eta_j=2^{E-j}$ and
$h_{ij}=\eta_jr_i$.

\paragraph{Algorithm 1 (one trial with grading $\theta$, cap $K$, accuracy
$\varepsilon$ and stage limit $J$).}
Start with $B^{(0)}=X$, center $c^{(0)}=\ell$, and $U$ the smallest objective
value of a feasible point known so far (at least $F(\ell)$). For
$j=0,1,\dots,J$:
\begin{enumerate}[label=(\roman*)]
\item for $i\in P$ let $G_i^{(j)}$ be the graded grid of $B_i^{(j)}$ with center
$c_i^{(j)}$, mesh $h_{ij}$ and grading $\theta$; for $i\notin P$ let
$G_i^{(j)}=\{\ell_i,u_i\}$; if some grid would exceed $K$ nodes, \emph{abort} the
trial before forming any table;
\item compute $\beta_j$, a minimizer $y^{(j)}$ and all min-marginals by
Lemma~\ref{lem:dp}, and set $U\leftarrow\min\{U,F(y^{(j)})\}$;
\item if $U-\beta_j\le\varepsilon$, stop with \emph{success}; otherwise filter
with threshold $U$ to obtain $B^{(j+1)}$ and set $c^{(j+1)}=y^{(j)}$.
\end{enumerate}
By Proposition~\ref{prop:filter} the centers stay feasible and inside the
current subbox, and by Theorem~\ref{thm:certificate} every stage yields a
valid path certificate.

\subsection{Contraction and localization}

\begin{lemma}[Stage invariants]\label{lem:inv}
Assume weighted growth with constant $\gamma$ and minimizer $x^*$, let
$\bar\kappa\ge\max\{1,1/\gamma\}$ with $8\bar\kappa\theta^2\le1$, and write
$a_j=n_P\eta_j^2$. At every stage $j$ reached by the trial:
\begin{enumerate}[label=(\roman*)]
\item $x^*\in B^{(j)}$ and $\beta_j\le\OPT$;
\item $\norm{c^{(j)}-x^*}_{\mathbf L}^2\le4\bar\kappa a_j$;
\item $\norm{y^{(j)}-x^*}_{\mathbf L}^2\le\bar\kappa a_j$;
\item $U-\beta_j\le F(y^{(j)})-\beta_j=D(y^{(j)})\le\frac9{16}a_j$;
\item every grid point $z$ with $Q(z)\le U$ (after the update at stage $j$)
satisfies $\norm{z-x^*}_{\mathbf L}^2\le\frac{17}{15}\bar\kappa a_j$.
\end{enumerate}
\end{lemma}

\begin{proof}
(i) follows from Proposition~\ref{prop:filter}, since $U\ge\OPT$, and from
Proposition~\ref{prop:cellwise}. Next, every grid point $z$ satisfies, with
$c=c^{(j)}$,
\begin{equation}\label{eq:energy}
D(z)\le\frac{a_j}4+\frac{\theta^2}2\Bigl(\norm{z-x^*}_{\mathbf L}^2
+\norm{c-x^*}_{\mathbf L}^2\Bigr).
\end{equation}
Indeed, for $i\in P$ Lemma~\ref{lem:graded}(a) and $L_ih_{ij}^2\le\eta_j^2$ give
$d_i(z_i)\le\frac{L_i}8(h_{ij}+\theta\abs{z_i-c_i})^2\le\frac{\eta_j^2}4
+\frac{\theta^2}4L_i(z_i-c_i)^2$, and
$(z_i-c_i)^2\le2(z_i-x_i^*)^2+2(c_i-x_i^*)^2$; for $i\notin P$, $d_i=0$.

(ii) For $j=0$, $\norm{\ell-x^*}_{\mathbf L}^2\le\sum_{i\in P}L_is_i^2\le
\sum_{i\in P}L_ir_i^2\eta_0^2\le a_0$. For $j\ge1$ it follows from (iii) at
stage $j-1$, because $a_{j-1}=4a_j$.

(iii) Write $Y=\norm{y^{(j)}-x^*}_{\mathbf L}^2$ and $a=a_j$. Since
$\theta^2/2\le1/(16\bar\kappa)\le\gamma/16$, growth, (i), $Q(y^{(j)})=\beta_j$,
\eqref{eq:energy} and (ii) give
$\gamma Y\le F(y^{(j)})-\OPT\le D(y^{(j)})\le\frac a4+\frac\gamma{16}(Y+4\bar\kappa a)$,
so $Y\le\frac4{15}(\gamma^{-1}+\bar\kappa)a\le\frac8{15}\bar\kappa a$.

(iv) By \eqref{eq:energy}, (ii), (iii) and $\theta^2\bar\kappa\le\frac18$,
$D(y^{(j)})\le\frac a4+\frac{\theta^2}2\cdot5\bar\kappa a\le\frac9{16}a$; and
$U\le F(y^{(j)})$.

(v) Let $Z=\norm{z-x^*}_{\mathbf L}^2$. Using $Q(z)\le U\le F(y^{(j)})$,
$\beta_j\le\OPT$, (iv) and \eqref{eq:energy},
$\gamma Z\le F(z)-\OPT=Q(z)+D(z)-\OPT\le D(y^{(j)})+D(z)\le\frac{13}{16}a
+\frac\gamma{16}(Z+4\bar\kappa a)$, so
$Z\le(\frac{13}{15}\gamma^{-1}+\frac4{15}\bar\kappa)a\le\frac{17}{15}\bar\kappa a$.
\end{proof}

The proof shows where filtering is essential. Part (v) bounds every grid point
whose corrected value can keep an interval alive, and the retained subbox
therefore shrinks at the same rate as the mesh.

\begin{lemma}[Localization and grid size]\label{lem:states}
Under the hypotheses of Lemma~\ref{lem:inv}, for every stage $j$ followed by
filtering and every $i\in P$,
\[
B_i^{(j+1)}\subseteq\bigl[y_i^{(j)}-R_{ij},\,y_i^{(j)}+R_{ij}\bigr],\qquad
R_{ij}=8\sqrt{\bar\kappa n_P}\,h_{ij}+[i\in I_Z].
\]
Consequently every grid of the trial has at most
\begin{equation}\label{eq:statecap}
K(\theta,n_P)=10\,\theta^{-1}\lceil\log_2(n_P+2)\rceil
\end{equation}
nodes per coordinate, independently of the stage and of $\varepsilon$.
\end{lemma}

\begin{proof}
Write $y=y^{(j)}$, $c=c^{(j)}$, $h=h_{ij}$ and $\rho=\sqrt{\bar\kappa n_P}$. Let
$[v,v']$ be a retained interval of coordinate $i$ with $m_i(v)\le U$, and let
the grid point $z$ attain $m_i(v)$, so $z_i=v$. Because $\sqrt{L_i}>1/(2r_i)$,
a bound $\sqrt{L_i}\abs\delta\le\lambda\rho\eta_j$ implies
$\abs\delta<2\lambda\rho h$; so Lemma~\ref{lem:inv}(ii), (iii), (v) give
$\abs{v-y_i}<2(1+\sqrt{17/15})\rho h$ and $\abs{v-c_i}<2(2+\sqrt{17/15})\rho h$.
An integer interval of length one lies within $\abs{v-y_i}+1$ of $y_i$.
Otherwise the interval has length at most $h+\theta\abs{v-c_i}$
(Lemma~\ref{lem:graded}(a)), and $\theta\sqrt{\bar\kappa}\le1/\sqrt8$, $h\le\rho h$
and $\sqrt{n_P}\le\rho$ give that each of its points is within
$\bigl(3+2\sqrt{17/15}+(2+\sqrt{17/15})/\sqrt2\bigr)\rho h<7.3\rho h$ of $y_i$.
The retained hull lies in the same interval.

At stage $j+1$ the grid of coordinate $i$ is centered at $y_i$ with mesh $h/2$,
and each side has length at most $R_{ij}$. With $H$ as in
Lemma~\ref{lem:graded}(b) for mesh $h/2$, $\theta R_{ij}/H\le\theta+16\theta\rho
\le\frac14+4\sqrt{2n_P}$, so the grid has at most
$1+2\lceil\frac4\theta\ln(\frac54+4\sqrt{2n_P})\rceil\le\theta^{-1}\varphi(n_P)$
nodes, with $\varphi(m)=\frac34+8\ln(\frac54+4\sqrt{2m})$, using $3\le3/(4\theta)$.
Now $\varphi(1)<16.3\le10\lceil\log_23\rceil$, $\varphi(2)<18.6\le10\log_24$, and
for $m\ge2$, $\varphi'(m)\le4/m$ while
$\frac{d}{dm}10\log_2(m+2)\ge7.2/m$; hence $\varphi(m)\le10\log_2(m+2)$. Stage
$0$ has at most three nodes per coordinate by Lemma~\ref{lem:graded}(c),
since $h_{i0}\ge s_i$.
\end{proof}

With a common mesh $h_j=s2^{-j}$ for all coordinates and the Euclidean
condition number $\kappa$, the same proof gives the radius
$4.2\sqrt{n\kappa}\,h_j$ and the cap $8\theta^{-1}\lceil\log_2(n+2)\rceil$;
the implementation of Section~\ref{sec:computation} uses this variant.
\subsection{Sharpness of the localization radius}
The slack $U-\beta$ can be as large as a sum of $n$ correction terms,
while a single coordinate interval is removed only when its own excess
exceeds the whole slack. The following family shows that the radius
$\Theta(\sqrt{n\kappa}\,h_j)$ of Lemmas~\ref{lem:states}
and~\ref{lem:commonmesh} is therefore attained, so that filtering alone
leaves at least of order $\sqrt{n\kappa}$ nodes per coordinate on uniform
grids; the graded grids reduce this to
$O(\sqrt{\bar\kappa}\log(n+2))$ (Lemma~\ref{lem:states}), which makes
Theorem~\ref{thm:approx} a fixed-parameter bound (Remark~\ref{rem:fpt}).

\begin{proposition}[Sharpness of localization]\label{prop:sharp}
Let $m\ge1$, $n=2m$, $\Lambda>0$ and $0<g\le\Lambda/2$, and on
$X=[-1,1]^n$, with all coordinates continuous and named
$(u_1,v_1,\dots,u_m,v_m)$, let
\[
F(x)=\sum_{k=1}^m\Bigl[\frac\Lambda2(u_k^2+v_k^2)+(\Lambda-2g)u_kv_k\Bigr].
\]
Then $x^*=0$, $\OPT=0$, the smallest valid curvature bounds are
$L_i=\Lambda$, the largest growth constant is $g$, and $\kappa=\Lambda/g$
takes every value in $[2,\infty)$. Let $G$ be a grid for a subbox that
contains $[-R_0,R_0]^n$, and let every $G_i$ contain $0$ with
$w_i(0)\ge h$. Then for every incumbent value $U\ge0$, every coordinate $i$
and every node $a\in G_i$ with
\[
|a|\le\min\Bigl\{R_0,\ h\sqrt{(n-2)\kappa/16}\Bigr\}
\]
we have $m_i(a)\le U$. Hence every grid interval with such an endpoint is
retained.
\end{proposition}

\begin{corollary}[Filtered uniform grids need order $\sqrt{n\kappa}$ nodes]
\label{cor:uniformgrid}
Take the instance of Proposition~\ref{prop:sharp} with $n\ge4$, and run the
stages of a trial with a common mesh (Lemma~\ref{lem:commonmesh}) without a
cap, with $h_j=s2^{-j}=2^{1-j}$, uniform grids (the construction of
Definition~\ref{def:graded} with $\theta=0$) and initial center $0$. Let
$k_0=\lfloor\sqrt{(n-2)\kappa/16}\rfloor$. At every stage the corrected
minimizer is $0$; at every stage $j\ge1$ with $(k_0+1)h_j\le1$ that ends by
filtering, the filtered box $X^{(j+1)}$ contains
$[-(k_0+1)h_j,(k_0+1)h_j]^n$, and the grid of stage $j+1$, if that stage is
run, has at least $4k_0+5\ge\sqrt{(n-2)\kappa}+1$ nodes in every coordinate.
With graded grids of grading $\theta\in(0,1/4]$ instead: if a stage $j$ ends
by filtering, its box $X^{(j)}$ contains $[-R,R]^n$, where
$R=h_j\sqrt{(n-2)\kappa/16}$, and its grids contain $0$ with
$w_i(0)\ge h_j$, then the grid of stage $j+1$, if that stage is run, has at
least $1+\ln(1+2\theta R/h_j)/\ln(1+\theta)$ nodes in every coordinate.
\end{corollary}

Proofs are in Appendix~\ref{app:growth}. Uniform grids also give
accuracy-independent counts without any grading condition: under point
growth, UC (Algorithm~\ref{alg:uc}), which filters unions of uniform cells,
has at most $12(2\sqrt{n\kappa}+1)$ nodes per coordinate at every stage
(Theorem~\ref{thm:cells} with $r=1$ and $\kappa_S=\kappa$). The resulting
running-time bound is polynomial in $I+q$ and $\kappa$ for fixed $p$, but
because of the factor $\sqrt n$ per coordinate it is not a fixed-parameter
bound.


\subsection{Unknown growth: capped trials}
The algorithm cannot use $8\bar\kappa\theta^2\le1$ directly because $\gamma$ is
unknown. It tries geometrically decreasing gradings and relies on the caps to
bound the work of unsuitable trials.

\paragraph{Algorithm 2 (CT$(\varepsilon)$, capped trials).}
If $P=\emptyset$, run one stage with $G_i=\{\ell_i,u_i\}$ for all $i$, which is
exact. Otherwise let $J$ be the least integer $J\ge0$ with
$\frac9{16}n_P\eta_0^2\,4^{-J}\le\varepsilon$, and for $\mu=2,3,\dots$ run a trial
with $\theta=2^{-\mu}$, cap $K_\mu=K(2^{-\mu},n_P)$ and stage limit $J$, until a
trial succeeds. Return the final incumbent $\hat x$, the bound $\beta$ of the
successful stage and the grids of the successful trial.

\begin{theorem}[Certified approximation]\label{thm:approx}
Let~\eqref{eq:model} have explicit rational quadratic factors given as
monomial lists, a tree decomposition of maximum bag size $p$, curvature bounds
$L_i=\max\{0,\partial_{ii}F\}$, and let $\varepsilon=2^{-q}$.
\begin{enumerate}[label=(\alph*)]
\item On every instance, CT$(\varepsilon)$ halts and returns a feasible
rational $\hat x$ and a valid path certificate with
$\beta\le\OPT\le F(\hat x)\le\beta+\varepsilon$.
\item If the instance has weighted growth with condition number $\bar\kappa$,
CT$(\varepsilon)$ succeeds at a trial with $2^\mu\le6\sqrt{\bar\kappa}$ and uses at
most
\[
f(p,\bar\kappa)\,(I+q+1)^5,\qquad
f(p,\kappa)=c_0\,(c_1p\sqrt\kappa)^p\,\kappa\,(1+\log_2\kappa)^2,
\]
bit operations, for absolute constants $c_0,c_1$. The certificate has size,
and is checked in time, within the same bound. Neither $\gamma$ nor $\bar\kappa$
is an input, and the bound holds in particular with $\kappa\ge\bar\kappa$.
\end{enumerate}
\end{theorem}

\begin{proof}
(a) Validity follows from Theorem~\ref{thm:certificate}. For termination,
take $\mu\ge2$ with $2^{-\mu}\le\min_{i\in P}h_{iJ}/s_i$. In this trial every
step at stage $j\le J$ is at least $h_{ij}$ (continuous) or
$\max\{1,\lfloor h_{ij}\rfloor\}$ (integer), so each grid has at most
$3+2/\theta<K_\mu$ nodes and the trial is not aborted; at stage $J$ every grid
interval has length at most $h_{iJ}+\theta s_i\le2h_{iJ}$, so
$U-\beta_J\le D(y^{(J)})\le\sum_{i\in P}L_ih_{iJ}^2/2\le\frac89\varepsilon$
(for integer coordinates with $h_{ij}<1$ the grid has at most $s_i+1<1/\theta+1$
nodes).

(b) Let $\mu^*=\max\{2,\lceil\log_4(8\bar\kappa)\rceil\}$, so $8\bar\kappa\theta^2\le1$
and $2^{\mu^*}\le6\sqrt{\bar\kappa}$. By Lemma~\ref{lem:states} trial $\mu^*$ is
never aborted, and by Lemma~\ref{lem:inv}(iv) it succeeds at stage $J$ at the
latest. Every trial $\mu\le\mu^*$ runs at most $J+1$ stages with at most $K_\mu$
nodes per coordinate, because an oversized grid is detected while it is
generated. By Lemma~\ref{lem:dp} this costs
$O\bigl((J+1)p(\abs{\mathcal A}+n+N)K_\mu^p\bigr)$ table operations, and
$\sum_{\mu\le\mu^*}K_\mu^p\le2K_{\mu^*}^p$. Since
$\sup_{t\ge0}t^pe^{-t}=(p/e)^p$,
\begin{equation}\label{eq:logabsorb}
\lceil\log_2(n_P+2)\rceil^p\le2p^p(n+2),
\end{equation}
so the table work is at most $c\,(J+1)\,p\,I^2(60\,p\sqrt{\bar\kappa})^p$ with
$J=O(I+q+1)$ and an absolute constant $c$. Appendix~\ref{app:growth} shows that all numbers have
denominators dividing $\Gamma_X2^\alpha$, where $\Gamma_X$ is the product of the
endpoint denominators and $\alpha=J+O(I)+\mu K_\mu$, so they do not grow with
the number of stages, and that every table entry is an integer of
$O((1+\mu2^\mu)(I+q+1))$ bits over a common denominator. Multiplying the counts
proves the bound.
\end{proof}

\begin{remark}[In what sense the result is fixed-parameter tractable]
\label{rem:fpt}
The running time is $f(p,\bar\kappa)(I+q+1)^5$ on every instance with weighted
growth, where $\bar\kappa$ is a property of the instance that the algorithm
neither receives nor computes; $f$ is computable and the exponent $5$, which
reflects schoolbook arithmetic, depends on neither parameter. The accuracy
enters only through $q=\log_2(1/\varepsilon)$. The term
$\lceil\log_2(n+2)\rceil^p$ in the table size is absorbed
by~\eqref{eq:logabsorb}, and no bound on the number of bags containing a
coordinate, on Hessian norms or on the number of local minima is used. If only
a treewidth bound is known, the decomposition of \cite{Korhonen2021} gives the
same result with $p=2\,\mathrm{tw}+2$. Section~\ref{sec:limits} shows that
neither parameter can be dropped and that the exponent of $\kappa$ must grow
linearly with $p$.
\end{remark}

\begin{remark}[Relation to the cluster problem]\label{rem:cluster}
In branch-and-bound with second-order convergent lower bounds, the number of
boxes that remain near a nondegenerate minimizer is independent of the
tolerance; this is the analysis of the cluster problem
\cite{DuKearfott1994,Neumaier2004,WechsungSchaberBarton2014,KannanBarton2017}.
That count is local, asymptotic and exponential in the dimension and in the
conditioning of the Hessian. Lemmas~\ref{lem:inv} and~\ref{lem:states} give a
global, non-asymptotic version for product grids on a tree decomposition. Because
filtering is coordinatewise, the exponential dependence falls on the bag size
$p$ rather than on $n$, and the condition $8\bar\kappa\theta^2\le1$ plays the role
of the bound on the second-order prefactor.
\end{remark}

\subsection{Examples}
The first example shows that the theorem covers nonconvex problems with many
local minima and treewidth larger than one.

\begin{example}[Coupled indefinite blocks]\label{ex:family}
For a block $b$ with variables $(u_b,v_b,\rho_b)\in[0,1]^3$ let
$\phi(u,v,\rho)=u^2+v^2-4uv+\frac14(u+v)+(\rho-\frac u2)^2$, and for a graph
$\Gamma$ on the blocks with maximum degree $\Delta$ let
\[
F=\sum_b\phi(u_b,v_b,\rho_b)+\frac1{16\Delta}\sum_{bb'\in\Gamma}(u_b-u_{b'})^2 .
\]
Then (a) $x^*=(1,1,\frac12)$ in every block is the unique minimizer and
$F(x)-\OPT\ge\frac12\norm{x-x^*}^2$; (b) $L\le\frac{21}8$, so
$\kappa\le\frac{21}4$; (c) each of the $2^m$ points with
$(u_b,v_b,\rho_b)\in\{(0,0,0),(1,1,\frac12)\}$ in all $m$ blocks is a strict
local minimizer; (d) a tree decomposition of $\Gamma$ of width $w$ yields one of
$F$ with $p=\max\{w+1,3\}$. Each block Hessian is indefinite and the minimizer
is not a vertex. For $\Gamma$ a $k\times m'$ grid graph, Theorem~\ref{thm:approx}
applies with $p=\max\{k+1,3\}$ and $\kappa\le\frac{21}4$, although the instance
has $2^{km'}$ strict local minima. The proof is in Appendix~\ref{app:growth}.
\end{example}

The minimizer of this family is easy to guess block by block; bounded $\kappa$
forces a flipped block to cost a fixed amount on average. The second example
separates the grids from exact parametric dynamic programming.

\begin{example}[Expanding-box chain]\label{ex:chain}
For $m\ge2$ consider continuous variables $S_t\in[0,2^t-1]$ and
$z_t\in[0,1]$, $t=1,\dots,m$, with $S_0=0$, and
\[
G_m(S,z)=S_m^2+\sum_{t=1}^m\bigl(S_t-2S_{t-1}-z_t\bigr)^2
+\frac18\sum_{t=1}^mz_t(1-z_t).
\]
The bags $\{S_{t-1},S_t,z_t\}$ form a path decomposition with $p=3$. The unique
minimizer is $0$, with growth constant $g=1/8$ and coordinate curvature $L=10$,
so $\kappa\le80$ for every $m$ (Proposition~\ref{lim:prop:messages}). By
Theorem~\ref{thm:approx}, a $2^{-q}$-approximation with certificate takes time
polynomial in $m$ and $q$. In contrast, the exact Bellman message for $S_m$ is a
piecewise quadratic function with at least $2^{m-1}$ pieces
(Proposition~\ref{lim:prop:messages}). Because $\bar\kappa$ is invariant under
rescaling, the same holds for the unit-box version $s_t=2^{-t}S_t$, whose
Euclidean condition number is $\Theta(4^m)$. Section~\ref{sec:computation}
reports the certified computation up to $m=64$.
\end{example}
